  \chapter{El subsistema $\RCA$}

\begin{definition}[$\RCA$]
	Se define al subsistema de $\ARITMETICA$ de \emph{axiomas de comprensión recursiva}, o $\RCA$, como el conjunto de axiomas:
	\[
		\RCA\coloneqq\PEANO\cup\INDUCCION{\Sigma_{1}^{0}}\cup\COMPRENSION{\Delta_{1}^{0}}
	\]
\end{definition}

\begin{theorem}
	Sea $S\subseteq\CPOTENCIA{\omega}$ una $\omega$-estructura. $S$ es un $\omega$-modelo de $\RCA$ si y solo si es un ideal de Turing.
\end{theorem}

\begin{proof}
	Para la necesidad debemos demostrar lo siguiente:

	\begin{enumerate}
		\item $S\neq\varnothing$

		Consideremos la fórmula $\varphi$ dada por $\neg\pa*{n = n}$. Notemos que $\varphi\in\Delta_{0}^{0}$ y luego $\varphi$ es $\Sigma_{1}^{0}$ y $\Pi_{1}^{0}$. Por lo tanto, por $\COMPRENSION{\Delta_{1}^{0}}$ sabemos que existe el conjunto $X\in S$ definido por $\varphi$. Se tiene entonces lo siguiente:
		\[
			X =\set*{n:\varphi\pa*{n}} =\set*{n: n\neq n} =\varnothing
		\]

		Por lo tanto, $\varnothing\in S$, y luego $S\neq\varnothing$.

		\item Para cada $X, Y\in S$ se satisface que $X\oplus Y\in S$.

		Consideremos a la fórmula $\varphi$ dada por:
		\[
			\pa*{\exists m < n + 1}\pa*{\pa*{m\in X\land n = 2m}\lor\pa*{m\in Y\land n = 2m + 1}}
		\]
		Notemos que $\varphi\in\Delta_{0}^{0}$, y entonces por $\COMPRENSION{\Delta_{1}^{0}}$ existe $Z\in S$ definido por $\varphi$. Se obtiene entonces lo siguiente:

		\begin{align*}
			Z =&\set*{n:\varphi\pa*{n}}\\
			=&\set*{n:\pa*{\exists m < n + 1}\pa*{\pa*{m\in X\land n = 2m}\lor\pa*{m\in Y\land n = 2m + 1}}}\\
			=&\set*{n:\pa*{\exists m < n + 1}\pa*{m\in X\land n = 2m}\lor\pa*{\exists m < n + 1}\pa*{m\in Y\land n = 2m + 1}}\\
			=&\set*{n:\pa*{\exists m < n + 1}\pa*{m\in X\land n = 2m}}\cup\set*{n:\pa*{\exists m < n + 1}\pa*{m\in Y\land n = 2m + 1}}\\
			=&\set*{2m: m\in X}\cup\set*{2m + 1: m\in Y}\\
			=& X\oplus Y
		\end{align*}

		Por lo tanto, $X\oplus Y\in S$.

		\item para cada $X, Y\in\CPOTENCIA{\omega}$, si $X\leq_{T} Y$ y $Y\in S$, entonces $X\in S$.

		Ya que $X\leq_{T} Y$, entonces por el teorema de Post (\ref{theo: teorema_de_post}) sabemos que existe una fórmula $\varphi$ con parámetro $Y$ la cual es $\Sigma_{1}^{0}$ y $\Pi_{1}^{0}$ y que define a $X$. Puesto que $Y\in S$, entonces por $\COMPRENSION{\Delta_{1}^{0}}$ se tiene que existe $Z\in S$ definida por $\varphi$. Por lo tanto, $X = Z\in S$.
	\end{enumerate}

	Probemos ahora la suficiencia. Observemos que puesto que $S$ es una $\omega$-estructura, entonces se satisfacen directamente $\PEANO$ y $\INDUCCION{\Sigma_{1}^{0}}$. Luego, es únicamente necesario demostrar que se satisface el esquema $\COMPRENSION{\Delta_{1}^{0}}$.

	Sea $\varphi\pa*{n}\in\Delta_{1}^{0}$ con $\mathrm{Free}\bra*{\varphi} =\set*{n}$, entonces demostremos que existe $X\in S$ definido por $\varphi$. Debido a que $S\neq\varnothing$, podemos suponer que $\varphi\pa*{n, Y_{1},\dotsc, Y_{k}}$ para algunos parámetros $Y_{1},\dotsc, Y_{k}\in S$ y $k\in\omega^{+}$. Consideremos al conjunto $X\in\CPOTENCIA{\omega}$ definido por $\varphi$, entonces por el teorema de Post generalizado (\ref{theo: teorema_de_post_generalizado}) sabemos que $X\leq_{T}Y_{1}\oplus\dotsb\oplus Y_{k}$. Ya que $Y_{1},\dotsc, Y_{k}\in S$, entonces $Y_{1}\oplus\dotsb\oplus Y_{k}\in S$, y por lo tanto, $X\in S$.
\end{proof}

\section{Teoría de números elemental}

De ahora en adelante, todos los siguientes teoremas serán demostrados dentro del subsistema $\RCA$, a menos que se diga lo contrario.

\begin{theorem}[Propiedades básicas de la aritmética]
	Se satisfacen las siguientes cerraduras universales.

	\begin{enumerate}
		\item $\pa*{n + m} + p = n +\pa*{m + p}$
		\item $0 + n = n$
		\item $n + 1 = 1 + n$
		\item $n + m = m + n$
		\item $n\pa*{m + p} = nm + np$
		\item $\pa*{nm}p = n\pa*{mp}$
		\item $\pa*{n + m}p = np + mp$
		\item $0\cdot n = 0$
		\item $1\cdot n = n$
		\item $nm = mn$
		\item $\pa*{n < m\land m < p}\to n < p$
		\item $n < m\leftrightarrow n + 1 < m + 1$
		\item $n\neq 0\to 0 < n$
		\item $\pa*{n < m}\lor\pa*{n = m}\lor\pa*{m < n}$
		\item $\neg\pa*{n < n}$
		\item $ n < m\leftrightarrow n + p < m + p$
		\item $n < n + m + 1$
		\item $n + p = m + p\to n = m$
		\item $\pa*{p\neq 0\land n < m}\to np < mp$
		\item $\pa*{p\neq 0\land np < mp}\to n < m$
		\item $\pa*{p\neq 0\land np = mp}\to n = m$
		\item $n < m\to\pa*{\exists p < m}\pa*{n + p + 1 = m}$
		\item $n\neq 0\to\pa*{\exists m < n}\pa*{m + 1 = n}$
	\end{enumerate}
\end{theorem}

\begin{proof}
	Cada una de estas afirmaciones se demostrarán por $\INDUCCION{\Sigma_{1}^{0}}$ notando que cada una de las fórmulas es $\Delta_{0}^{0}$, y luego $\Sigma_{1}^{0}$. Además, siempre haremos referencia a los axiomas básicos de la definición \ref{def: axiomas_base}.
	\begin{enumerate}
		\item $\pa*{n + m} + p = n +\pa*{m + p}$
		
		Mostermos por inducción sobre $p$. Para el caso base $p = 0$ se satisface lo siguiente por el axioma (\ref{def_axiomas: el_cero_es_identidad}):
		\[
			\pa*{n + m} + 0 = n + m = n +\pa*{m + 0}
		\]

		Ahora, sea $p$ tal que $\pa*{n + m} + p = n +\pa*{m + p}$, y mostremos que se satisface para $p + 1$. Por el axioma (\ref{def_axiomas: sucesor_es_asociativo}) y la hipótesis de inducción se sigue:
		\begin{align*}
			\pa*{n + m} + \pa*{p + 1} &= \pa*{\pa*{n + m} + p} + 1 \\
			&= \pa*{n +\pa*{m + p}} + 1 \\
			&= n +\pa*{\pa*{m + p} + 1}\\
			&= n +\pa*{m +\pa*{p + 1}}
		\end{align*}

		Se sigue por $\INDUCCION{\Sigma_{1}^{0}}$ que para cada $n, m, p$ se satisface $\pa*{n + m} + p = n +\pa*{m + p}$.

		\item $0 + n = n$
		
		Mostremos por inducción sobre $n$. Para el caso base $n = 0$ se sigue que $0 + 0 = 0$ por el axioma (\ref{def_axiomas: el_cero_es_identidad}).

		Sea ahora $n$ tal que $0 + n = n$, y mostremos que se satisface para $n + 1$. Por el axioma (\ref{def_axiomas: sucesor_es_asociativo})
		\[
			0 +\pa*{n + 1} =\pa*{0 + n} + 1 = n + 1
		\]

		Se sigue por $\INDUCCION{\Sigma_{1}^{0}}$ que para cada $n$ se satisface $0 + n = n$.

		\item $n + 1 = 1 + n$
		
		Mostremos por inducción sobre $n$. Para el caso base $n = 0$ se tiene por el punto anterior y por el axioma (\ref{def_axiomas: el_cero_es_identidad}) que $0 + 1 = 1 = 1 + 0$.

		Sea ahora $n$ tal que $n + 1 = 1 + n$ y mostremos que se satisface para $n + 1$. Por el axioma (\ref{def_axiomas: sucesor_es_asociativo}) y la hipótesis de inducción se satisface lo siguiente:
		\[
			\pa*{n + 1} + 1 = \pa*{1 + n} + 1 = 1 + \pa*{n + 1}
		\]

		Se sigue por $\INDUCCION{\Sigma_{1}^{0}}$ que para cada $n$ se satisface $n + 1 = 1 + n$.

		\item $n + m = m + n$
		
		Mostremos por inducción sobre $n$. Para el caso base $n = 0$ se tiene por el (2) y por el axioma (\ref{def_axiomas: el_cero_es_identidad}) que $0 + m = m = m + 0$.

		Sea ahora $n$ tal que $n + m = m + n$ y mostremos que se satisface para $n + 1$. Por el punto anterior, el primer punto y la hipótesis de inducción se satisface lo siguiente:
		\[
			\pa*{n + 1} + m = n +\pa*{1 + m} = n +\pa*{m + 1} =\pa*{n + m} + 1 =\pa*{m + n} + 1 = m +\pa*{n + 1}
		\]

		Se sigue por $\INDUCCION{\Sigma_{1}^{0}}$ que para cada $n$ se satisface $n + m = m + n$.

		\item $n\pa*{m + p} = nm + np$
		
		Mostremos por inducción sobre $p$. Para el caso base $p = 0$ se tiene por los axiomas (\ref{def_axiomas: cero_es_absorbente}) y (\ref{def_axiomas: el_cero_es_identidad}) lo siguiente:
		\[
			n\pa*{m + 0} = nm = nm + 0 = nm + n\cdot 0
		\]

		Sea ahora $p$ tal que $n\pa*{m + p} = nm + np$, y mostremos que se satisface para $p + 1$. Por los axiomas (\ref{def_axiomas: sucesor_es_asociativo}), (\ref{def_axiomas: cero_es_absorbente}) y (\ref{def_axiomas: producto_es_distributivo}), y la hipótesis de inducción se tiene lo siguiente:
		\begin{align*}
			n\pa*{m +\pa*{p + 1}} &= n\pa*{\pa*{m + p} + 1} \\
			&= n\pa*{m + p} + n \\
			&=\pa*{nm + np} + n \\
			&= nm +\pa*{np + n} \\
			&= nm + n\pa*{p + 1}
		\end{align*}

		Se sigue por $\INDUCCION{\Sigma_{1}^{0}}$ que para cada $n, m, p$ se satisface $n\pa*{m + p} = nm + np$.

		\item $\pa*{nm}p = n\pa*{mp}$
		
		Mostremos por inducción sobre $p$. Para el caso base $p = 0$ se tiene por el axioma (\ref{def_axiomas: cero_es_absorbente}) lo siguiente:
		\[
			\pa*{nm}\cdot 0 = 0 = n\cdot 0 = n\pa*{m\cdot 0}
		\]

		Sea ahora $p$ tal que $\pa*{nm}p = n\pa*{mp}$, y mostremos que se satisface para $p + 1$. Por el axioma (\ref{def_axiomas: producto_es_distributivo}), el punto anterior, y la hipótesis de inducción se tiene lo siguiente:
		\begin{align*}
			\pa*{nm}\pa*{p + 1} = \pa*{nm}p + nm = n\pa*{mp} + nm = n\pa*{mp + m} = n\pa*{m\pa*{p + 1}}
		\end{align*}

		Se sigue por $\INDUCCION{\Sigma_{1}^{0}}$ que para cada $n, m, p$ se satisface $\pa*{nm}p = n\pa*{mp}$.

		\item $\pa*{n + m}p = np + mp$
		
		Mostremos por inducción sobre $p$. Para el caso base $p = 0$ se tiene por los axiomas (\ref{def_axiomas: el_cero_es_identidad}) y (\ref{def_axiomas: cero_es_absorbente}) lo siguiente:
		\[
			\pa*{n + m}\cdot 0 = 0 = 0 + 0 = n\cdot 0 + m\cdot 0
		\]

		Sea ahora $p$ tal que $\pa*{n + m}p = np + mp$, y mostremos que se satisface para $p + 1$. Por el axioma (\ref{def_axiomas: producto_es_distributivo}), el primer y cuarto punto, y la hipótesis de inducción se tiene lo siguiente:
		\begin{align*}
			\pa*{n + m}\pa*{p + 1} &= \pa*{n + m}p + \pa*{n + m} \\
			&=\pa*{np + mp} +\pa*{n + m}\\
			&=\pa*{\pa*{np + mp} + n} + m\\
			&=\pa*{np +\pa*{mp + n}} + m\\
			&=\pa*{np +\pa*{n + mp}} + m\\
			&=\pa*{\pa*{np + n} + mp} + m\\
			&=\pa*{n\pa*{p + 1} + mp} + m\\
			&= n\pa*{p + 1} +\pa*{mp + m}\\
			&=n\pa*{p + 1} + m\pa*{p + 1}
		\end{align*}

		Se sigue por $\INDUCCION{\Sigma_{1}^{0}}$ que para cada $n, m, p$ se satisface $\pa*{n + m}p = np + mp$.

		\item $0\cdot n = 0$
		
		Mostremos por inducción sobre $n$. Para el caso base $n = 0$ se tiene por el axioma (\ref{def_axiomas: cero_es_absorbente}) que $0\cdot 0 = 0$.

		Sea ahora $n$ tal que $0\cdot n = 0$, y mostremos que se satisface para $n + 1$. Por los axiomas (\ref{def_axiomas: el_cero_es_identidad}) y (\ref{def_axiomas: producto_es_distributivo}), y la hipótesis de inducción se tiene lo siguiente:
		\[
			0\cdot\pa*{n + 1} = 0\cdot n + 0 = 0 + 0 = 0
		\]

		Se sigue por $\INDUCCION{\Sigma_{1}^{0}}$ que para cada $n$ se satisface $0\cdot n = 0$.

		\item $1\cdot n = n$
		
		Mostremos por inducción sobre $n$. Para el caso base $n = 0$ se tiene por el axioma (\ref{def_axiomas: cero_es_absorbente}) que $1\cdot 0 = 0$.

		Sea ahora $n$ tal que $1\cdot n = n$, y mostremos que se satisface para $n + 1$. Por el axioma (\ref{def_axiomas: producto_es_distributivo}) y la hipótesis de inducción se tiene lo siguiente:
		\[
			1\cdot\pa*{n + 1} = 1\cdot n + 1 = n + 1
		\]

		Se sigue por $\INDUCCION{\Sigma_{1}^{0}}$ que para cada $n$ se satisface $1\cdot n = n$.

		\item $nm = mn$
		
		Mostremos por inducción sobre $n$. Para el caso base $n = 0$ se tiene por el axioma (\ref{def_axiomas: cero_es_absorbente}) y por el octavo punto que $0\cdot m = 0 = m\cdot 0$.

		Sea ahora $n$ tal que $nm = mn$, y mostremos que se satisface para $n + 1$. Por el axioma (\ref{def_axiomas: producto_es_distributivo}), los puntos séptimo y noveno, y la hipótesis de inducción se tiene lo siguiente:
		\[
			\pa*{n + 1}m = nm + 1\cdot m = nm + m = mn + n = m\pa*{n + 1}
		\]

		Se sigue por $\INDUCCION{\Sigma_{1}^{0}}$ que para cada $n, m$ se satisface $nm = mn$.

		\item $\pa*{n < m\land m < p}\to n < p$
		
		Mostremos por inducción sobre $p$. Para el caso base $p = 0$ se tiene por vacuidad debido al axioma (\ref{def_axiomas: no_hay_negativos}) que $\pa*{n < m\land m < 0}\to n < 0$.

		Sea $p$ tal que $n < m$ y $m < p$ implican $n < p$, y mostremos para $p + 1$. Si suponemos que $n < m$ y $m < p + 1$, entonces por el axioma (\ref{def_axiomas: el_menor_o_igual}) se tienen dos casos: $m < p$ o $m = p$. Si $m < p$, por hipótesis de inducción se sigue que $n < p$; y por el mismo axioma (\ref{def_axiomas: el_menor_o_igual}) se tiene que $n < p + 1$. Si $m = p$, entonces $n < p$ por sustitución, de lo cual se implica $n < p + 1$ por el axioma (\ref{def_axiomas: el_menor_o_igual}). En cualquier caso se concluye que si $n < m$ y $m < p + 1$, entonces $n < p + 1$

		Se sigue por $\INDUCCION{\Sigma_{1}^{0}}$ que para cada $n, m, p$ se satisface que $n < m$ y $n < p$ implican $n < p$.

		\item $n < m\leftrightarrow n + 1 < m + 1$
		
		Mostremos por inducción sobre $m$. Para el caso base $m = 0$, la primer implicación ($n < 0\to n + 1 < 0 + 1$) se satisface por vacuidad por el axioma \ref{def_axiomas: no_hay_negativos}; si ahora suponemos que $n + 1 < 0 + 1$, entonces $n + 1 < 0$ o $n + 1 = 0$ por el axioma \ref{def_axiomas: el_menor_o_igual}. Si $n + 1 < 0$, entonces se satisface que $n < m$ por vacuidad debido al axioma (\ref{def_axiomas: no_hay_negativos}). Ahora, supongamos que $n + 1 = 0$ es verdadero, entonces por el axioma (\ref{def_axiomas: el_menor_o_igual}) se satisface que $n < n + 1 = 0$; lo cual es una contradicción por el axioma (\ref{def_axiomas: no_hay_negativos}), y luego se satisface que $n < m$ por vacuidad.

		Sea $m$ tal que $n < m$ si y solo si $n + 1 < m + 1$, y mostremos que se satisface para $m + 1$. Si suponemos primero que $n < m + 1$, entonces por el axioma (\ref{def_axiomas: el_menor_o_igual}) se satisface que $n < m$ o $n = m$. En el primer caso, por hipótesis de inducción, se satisface que $n + 1 < m + 1$, y por el mismo axioma y el punto anterior se satisface que $n + 1 < m + 1 <\pa*{m + 1} + 1$, es decir, $n + 1 <\pa*{m + 1} + 1$. En el segundo caso, por el axioma (\ref{def_axiomas: el_menor_o_igual}) sabemos que $m + 1 <\pa*{m + 1} + 1$, y puesto que $n = m$, entonces reemplazando se satisface que $n + 1 <\pa*{m + 1} + 1$. En cualquier caso se satisface que $n < m + 1$ implica $n + 1 <\pa*{m + 1} + 1$.

		Supongamos ahora que $n + 1 <\pa*{m + 1} + 1$, entonces por el axioma (\ref{def_axiomas: el_menor_o_igual}) se satisface que $n + 1 < m + 1$ o $n + 1 = m + 1$. En el primer caso, por hipótesis de inducción, se satisface que $n < m$, y puesto que $m < m + 1$ por el mismo axioma, se implica por el punto anterior que $n < m + 1$. Ahora, en el segundo caso, por el axioma (\ref{def_axiomas: sucesor_es_inyectiva}) se satisface que $n = m$, y puesto que $m < m + 1$, entonces reemplazando se cumple que $n < m + 1$. En cualquier caso se satisface que $n + 1 <\pa*{m + 1} + 1$ implica $n < m + 1$.

		Por lo tanto, se obtiene que $n < m + 1$ si y solo si $n + 1 <\pa*{m + 1} + 1$. Se sigue entonces por $\INDUCCION{\Sigma_{1}^{0}}$ que para cada $n, m$ se satisface que $n < m$ si y solo si $n + 1 < m + 1$.

		\item $n\neq 0\to 0 < n$
		
		Mostremos por inducción sobre $n$. Para el caso base se satisface $0\neq 0\to 0 < 0$ por vacuidad.

		Sea $n$ tal que $n\neq 0$ implica $0 < n$, y mostremos que se cumple para $n + 1$. Si $n + 1\neq 0$, entonces tenemos dos casos: $n = 0$ o $n\neq 0$. Si $n = 0$, entonces por el axioma (\ref{def_axiomas: el_menor_o_igual}) se tiene que $n = 0 < 0 + 1 = n + 1$. Si suponemos ahora el segundo caso, entonces por hipótesis de inducción la afirmación $n\neq 0$ implica $0 < n$. Por el mismo axioma (\ref{def_axiomas: el_menor_o_igual}) se tiene que $n < n + 1$, y por transitividad del décimoprimer punto se cumple que $0 < n + 1$. En cualquier caso $n + 1\neq 0$ implica $0 < n + 1$.

		Se sigue por $\INDUCCION{\Sigma_{1}^{0}}$ que para cada $n$ se satisface que $n \neq 0$ si y solo si $0 < n$.

		\item $\pa*{n < m}\lor\pa*{n = m}\lor\pa*{m < n}$
		\item $\neg\pa*{n < n}$
		\item $ n < m\leftrightarrow n + p < m + p$
		\item $n < n + m + 1$
		\item $n + p = m + p\to n = m$
		\item $\pa*{p\neq 0\land n < m}\to np < mp$
		\item $\pa*{p\neq 0\land np < mp}\to n < m$
		\item $\pa*{p\neq 0\land np = mp}\to n = m$
		\item $n < m\to\pa*{\exists p < m}\pa*{n + p + 1 = m}$
		\item $n\neq 0\to\pa*{\exists m < n}\pa*{m + 1 = n}$
	\end{enumerate}
\end{proof}

\comment{Demostraciones}{Demostrar las propiedades del teorema anterior.}


\begin{definition}[Mínimo]
	Sea $X$ un conjunto. Se dice que $X$ tiene un \emph{elemento mínimo} si se satisface la siguiente fórmula:
	\[
		\exists k\forall m\pa*{k\in X\land\pa*{m\in X\to k\leq m}}
	\]
\end{definition}

\begin{theorem}[Principio del buen orden]
	Todo conjunto $X$ no vacío, es decir tal que se satisface $\exists n\pa*{n\in X}$, tiene un elemento mínimo.
\end{theorem}

\begin{proof}
	Demostremos por contradicción. Consideremos a $X$ un conjunto no vacío, es decir tal que se satisface $\exists n\pa*{n\in X}$, para el cual se cumple lo siguiente:
	\begin{equation}
		\label{eq: buen_orden}
		\forall k\exists m\pa*{k\notin X\lor\pa*{m\in X\land m < k}}
	\end{equation}
	Sea $\varphi\pa*{p}$ la fórmula dada por $\pa*{\forall k\leq p}\pa*{k\notin X}$, la cual es $\Delta_{0}^{0}$. Mostremos por $\INDUCCION{\Sigma_{1}^{0}}$ que se satisface la fórmula $\forall p\varphi\pa*{p}$.

	\begin{itemize}
		\item (Base) Si $p = 0$, entonces debemos mostrar que $0\notin X$. Supongamos por el contrario que $0\in X$, entonces por la fórmula \eqref{eq: buen_orden} se sigue que existe $m$ tal que $m < 0$ y $m\in X$. Puesto que por los axiomas $\PEANO$ se tiene que para cada $n$ se satisface que $\neg\pa*{n < 0}$, entonces llegamos a una contradicción. Por lo tanto, necesariamente $0\notin X$.

		\item (Paso inductivo) Sea $n$ tal que se satisface $\varphi\pa*{n}$, entonces demostremos que se satisface $\varphi\pa*{n + 1}$.

		Supongamos que $n + 1\in X$, entonces por la fórmula \eqref{eq: buen_orden} se sigue que existe $m$ tal que $m\in X$ y $m < n + 1$, es decir, $m\leq n$. Ya que se satisface $\varphi\pa*{n}$ por hipótesis, se tiene entonces que $m\notin X$, lo cual contradice el hecho que $m\in X$ como se dedujo anteriormente. Por lo tanto, necesariamente $n + 1\notin X$.
	\end{itemize}

	Se sigue entonces por $\INDUCCION{\Sigma_{1}^{0}}$ que:
	\[
		\forall p\varphi\pa*{p}\equiv\forall p\pa*{\forall k\leq p}\pa*{k\notin X}\equiv\forall p\pa*{p\notin X}
	\]
	lo cual contradice el hecho que $X$ era no vacío.

	En conclusión, por contradicción, se prueba el teorema.
\end{proof}

\begin{definition}[Cota superior]
	Sea $X$ un conjunto. $X$ se dice estar \emph{acotado superiormente} si se satisface la siguiente fórmula:
	\[
		\exists n\forall m\pa*{m\in X\to m\leq n}
	\]
\end{definition}

\begin{definition}[Máximo]
	Sea $X$ un conjunto. Se dice que $X$ tiene un \emph{elemento máximo} si se satisface la siguiente fórmula:
	\[
		\exists k\forall m\pa*{k\in X\land\pa*{m\in X\to m\leq k}}
	\]
\end{definition}

\begin{theorem}[Principio del máximo]
	Todo conjunto $X$ no vacío y acotado superiormente tiene un elemento máximo.
\end{theorem}

\begin{proof}
	Defínase la fórmula $\varphi\pa*{n}$ dada por:
	\begin{equation}
		\label{eq: principio_maximo}
		\pa*{\forall k\leq n}\pa*{k\notin X}\lor\pa*{\exists p\leq n}\pa*{\forall k\leq n}\pa*{k\in X\to k\leq p}
	\end{equation}
	la cual es $\Delta_{0}^{0}$. Observemos que esta fórmula se interpreta como que dado $n$ se satisface que todos los elementos menores a $n$ no pertenecen a $X$ o el conjunto $X\cap\set{0,\dotsc, n}$ tiene máximo. Probemos entonces por $\INDUCCION{\Sigma_{1}^{0}}$ que se satisface la fórmula $\forall n\varphi\pa*{n}$.

	\begin{enumerate}
		\item (Base) Para $n = 0$ se tienen dos casos: $0\in X$ o $0\notin X$. Si se cumple $0\notin X$, entonces directamente se satisface la primera parte de la fórmula \eqref{eq: principio_maximo}. Si se satisface $0\in X$, entonces tomando $p = 0$ se satisface la segunda parte de la fórmula \eqref{eq: principio_maximo}. En cualquier caso obtenemos que se satisface $\varphi\pa*{0}$.

		\item (Paso inductivo) Sea $n$ tal que $\varphi\pa*{n}$. Entonces demostremos que se satisface que $\varphi\pa*{n + 1}$.

		Ya que se satisface $\varphi\pa*{n}$, entonces tenemos tres casos:

		\begin{itemize}
			\item Si se satisface $n + 1\in X$, entonces se satisface la segunda parte de la fórmula $\varphi\pa*{n + 1}$ tomando $p = n + 1$.

			\item Si se satisface $n + 1\notin X$ y para cada $k$ con $k\leq n$ se cumple que $k\notin X$, entonces claramente se cumple la primer parte de la fórmula $\varphi\pa*{n + 1}$.

			\item Si se satisface $n + 1\notin X$ y existe un elemento máximo de los elementos de $X$ acotados por $n$, entonces directamente se satisface la segunda parte de la fórmula $\varphi\pa*{n + 1}$.
		\end{itemize}

		Se sigue entonces que en cualquier caso se satisface $\varphi\pa*{n + 1}$. Luego, por $\INDUCCION{\Sigma_{1}^{0}}$ se cumple la fórmula $\forall n\varphi\pa*{n}$.

		Ya que $X$ es acotado superiormente, se deduce que existe $m$ tal que para cada $k$, si $k\in X$, entonces $k\leq m$. Por lo tanto, se satisface particularmente $\varphi\pa*{m}$.

		Ahora, puesto que $X$ es no vacío, entonces no se satisface la primera parte de $\varphi\pa*{m}$, de lo cual se sigue que existe $p$ con $p\leq m$ tal que es máximo de los elementos de $X$ acotados por $m$. Como $m$ ya es cota superior de $X$, entonces tal $p$ es un máximo de $X$.
	\end{enumerate}
\end{proof}

\section{Sistemas numéricos}
\subsection{Los números enteros}
En esta sección desarrollaremos la construcción de los números enteros, racionales y reales que son necesarios para poder abordar el estudio de los teoremas del Análisis Real que analizaremos en cada uno de los subsistemas.

Fuera de $\RCA$, recordemos que la característica más importante de los números enteros es que existe la resta, y que para cada número entero $m$ existen dos números naturales $n, k$ tales que $m = n - k$. Con esta propiedad podemos codificar cada número entero con tal par de números naturales que se interpretan de la forma anterior. Por ejemplo, si tenemos el par $\pa*{1, 2}$, entonces lo interpretamos como el número entero $1 - 2 = -1$. Sin embargo, notemos que existen varios pares que nos darían el mismo número entero en ese sentido. Para el caso del $-1$ se tiene que los pares $\pa*{1, 2}$, $\pa*{2, 3}$, $\pa*{3, 4}$ y $\pa*{4, 5}$ se interpretarían de la misma forma. En ese caso, decimos que todos estos pares son \emph{equivalentes} para representar al número $-1$. En general, para cada par de pares $\pa*{n, m}$ y $\pa*{p, q}$, estos son equivalentes si se interpretan como el mismo número entero, es decir, si $n - m = p - q$, o equivalentemente $n + q = p + m$. Con esta condición podemos luego definir una relación de equivalencia que nos codifique la \emph{igualdad} entre números enteros en $\RCA$.

\begin{definition}[Igualdad en $\ZZ$]
	Se define la relación $=_{\ZZ}\subseteq\NN^{2}\times\NN^{2}$ como sigue para cada $n, m, p, q\in\NN$:
	\[
		\inner*{n, m} =_{\ZZ}\inner*{p, q}\iff n + q = p + m
	\]
\end{definition}

Notemos que esta relación puede definirse por medio de una fórmula $\Delta_{0}^{0}$, por lo que existe debido a $\COMPRENSION{\Delta_{1}^{0}}$.

Ya que vimos que existen elementos diferentes que son \emph{iguales} mediante esta relación, entonces es necesario demostrar que esta relación es de equivalencia para que toda la construcción sea coherente.

\begin{theorem}
	\label{theo: relacion_equivalencia_z}
	La relación $=_{\ZZ}$ es de equivalencia.
\end{theorem}

\begin{proof}
	Demostremos entonces que $=_{\ZZ}$ es reflexiva, simétrica y transitiva. Consideremos a $n, m, p, q, r, s\in\NN$.

	\begin{enumerate}
		\item Reflexividad

		Puesto que $n + m = m + n$, entonces se sigue directamente $\inner*{n, m} =_{\ZZ}\inner*{n, m}$. Por lo tanto, $=_{\ZZ}$ es reflexiva.

		\item Simetría

		Supongamos que $\inner*{n, m} =_{\ZZ}\inner*{p, q}$, entonces por definición se satisface que $n + q = p + m$, o equivalentemente, $p + m = n + q$. De lo anterior se sigue que $\inner*{p, q} =_{\ZZ}\inner*{n, m}$. Por lo tanto, $=_{\ZZ}$ es simétrica.

		\item Transitividad

		Supongamos que $\inner*{n, m} =_{\ZZ}\inner*{p, q}$ y $\inner*{p, q} =_{\ZZ}\inner*{r, s}$, entonces se satisface que $n + q = p + m$ y $p + s = r + q$. Sumando ambas igualdades obtenemos lo siguiente:
		\[
			\begin{NiceArray}{c}
				\pa*{n + q} +\pa*{p + s} =\pa*{p + m} +\pa*{r + q}\\
				\pa*{n + s} +\pa*{p + q} =\pa*{r + m} +\pa*{p + q}\\
				n + s = r + m
			\end{NiceArray}
		\]
		Lo cual implica que $\inner*{n, m} =_{\ZZ}\inner*{r, s}$. Por lo tanto, $=_{\ZZ}$ es transitiva
	\end{enumerate}

	En conclusión, la relación $=_{\ZZ}$ es de equivalencia
\end{proof}

\begin{notation}
	Sean $A\subseteq\NN$ y $R\subseteq A^{2}$ de equivalencia, entonces para cada $x\in A$ denotaremos por $\bra*{x}_{R}$ a la clase de equivalencia de $x$, es decir:
	\[
		\bra*{x}_{R} =\set*{y\in A: xRy}
	\]
\end{notation}

Todo esto nos implica que dado cualquier $x\in\NN^{2}$, su clase de equivalencia $\bra*{x}_{=_{\ZZ}}$ es el conjunto de todos los pares de naturales que se interpretan como el mismo número entero. Más aún, puesto que $=_{\ZZ}$ es un conjunto que existe en $\RCA$, entonces la clase de equivalencia de cualquier par de naturales existe ya que puede definirse fácilmente por una fórmula $\Delta_{0}^{0}$.

Ya que tenemos el concepto de igualdad en los enteros, podemos definir lo que es la suma y producto en ellos. Fuera de $\RCA$, consideremos dos números enteros $x = n - m$ y $y = p - q$, entonces su suma y producto satisface lo siguiente:
\[
	\begin{NiceArray}{c}
		x + y =\pa*{n - m} +\pa*{p - q} =\pa*{n + p} -\pa*{m + q}\\
		xy =\pa*{n - m}\pa*{p - q} = np - nq - mp + mq =\pa*{np + mq} -\pa*{nq + mp}
	\end{NiceArray}
\]

De esta forma, podríamos definir la suma y producto de enteros como los pares $\pa*{n + p, m + q}$ y $\pa*{np + mq, nq + mp}$ respectivamente. Con esto definimos sus equivalentes en $\RCA$.

\begin{definition}
	\label{def: suma_producto_z_aux}
	Definimos las siguientes funciones:
	\[
		\begin{NiceArray}{cc}
			f:\NN^{2}\times\NN^{2}\to\NN & g:\NN^{2}\times\NN^{2}\to\NN\\
			\pa*{\inner*{n, m},\inner*{p, q}}\mapsto\inner*{n + p, m + q} &\pa*{\inner*{n, m},\inner*{p, q}}\mapsto\inner*{np + mq, nq + mp}
		\end{NiceArray}
	\]
\end{definition}

Por el momento no les denominamos suma ni producto debido a que necesitaremos redefinirlas cuando establezcamos propiamente el conjunto de números enteros.

En cambio, demostremos que estas dos operaciones se ``comportan bien'', es decir, que son independientes del representante, son conmutativas, asociativas y distributivas.

\begin{lemma}
	\label{lemma: independencia_representantes_z}
	Para cada $x, y, z, w\in\NN^{2}]$ se satisface que si $x =_{\ZZ} z$ y $y =_{\ZZ} w$, entonces $f\pa*{x, y} =_{\ZZ} f\pa*{z, w}$ y $g\pa*{x, y} =_{\ZZ} g\pa*{z, w}$.
\end{lemma}

\begin{proof}
	Sean $n, m, p, q, r, s, t, u\in\NN$ tales que $x =\inner*{n, m}$, $y =\inner*{p, q}$, $z =\inner*{r, s}$ y $w =\inner*{t, u}$.

	\begin{enumerate}
		\item Mostremos que $f\pa*{x, y} =_{\ZZ} f\pa*{z, w}$.

		Puesto que $x =_{\ZZ} z$ y $y =_{\ZZ} w$, entonces $n + s = r + m$ y $p + u = t + q$. Luego:
		\begin{align*}
			\pa*{n + s} +\pa*{p + u} &=\pa*{r + m} +\pa*{t + q}\\
			\pa*{n + p} +\pa*{s + u} &=\pa*{r + t} +\pa*{m + q}\\
			\inner*{n + p, m + q} &=_{\ZZ}\inner*{r + t, s + u}\\
			f\pa*{x, y} &=_{\ZZ} f\pa*{z, w}
		\end{align*}

		\item Mostremos que $g\pa*{x, y} =_{\ZZ} g\pa*{z, w}$.

		Puesto que $x =_{\ZZ} z$ y $y =_{\ZZ} w$, entonces $n + s = r + m$ y $p + u = t + q$. Luego:
		\begin{align*}
			p + u &= t + q\\
			n\pa*{p + u} &= n\pa*{t + q}\\
			np + nu &= nt + nq\\
			np + nu + mq &= nt + nq + mq\\
			np + mq + nu &= nq + mq + nt\\
			np + mq + nu + su &= nq + mq + nt + su\\
			np + mq +\pa*{n + s}u &= nq + mq + nt + su\\
			np + mq +\pa*{r + m}u &= nq + mq + nt + su\\
			np + mq + ru + mu &= nq + mq + nt + su\\
			np + mq + ru + mu + mt &= nq + mq + mt + nt + su\\
			np + mq + ru + mu + mt &= nq + m\pa*{t + q} + nt + su\\
			np + mq + ru + mu + mt &= nq + m\pa*{p + u} + nt + su\\
			np + mq + ru + mu + mt &= nq + mp + mu + nt + su\\
			np + mq + ru + mt &= nq + mp + nt + su\\
			np + mq + ru + mt + st &= nq + mp + nt + st + su\\
			np + mq + ru + mt + st &= nq + mp +\pa*{n + s}t + su\\
			np + mq + ru + mt + st &= nq + mp +\pa*{r + m}t + su\\
			np + mq + ru + mt + st &= nq + mp + rt + mt + su\\
			np + mq + ru + st &= nq + mp + rt + su\\
			\pa*{np + mq} +\pa*{ru + st} &=\pa*{nq + mp} +\pa*{rt + su}\\
			\inner*{np + mq, nq + mp} &=_{\ZZ}\inner*{rt + su, ru + st}\\
			g\pa*{x, y} &=_{\ZZ} g\pa*{z, w}
		\end{align*}
	\end{enumerate}
\end{proof}

\begin{lemma}
	\label{lemma: prop_suma_prod_z}
	Para cada $x, y, z\in\NN^{2}$ se satisface lo siguiente:

	\begin{enumerate}
		\item $f\pa*{x, y} = f\pa*{y, x}$
		\item $g\pa*{x, y} = g\pa*{y, x}$
		\item $f\pa*{x, f\pa*{y, z}} = f\pa*{f\pa*{x, y}, z}$
		\item $g\pa*{x, g\pa*{y, z}} = g\pa*{g\pa*{x, y}, z}$
		\item $g\pa*{x, f\pa*{y, z}} = f\pa*{g\pa*{x, y}, g\pa*{x, z}}$
	\end{enumerate}
\end{lemma}

\begin{proof}
	Sean $n, m, p, q, r, s\in\NN$ tales que $x =\inner*{n, m}$, $y =\inner*{p, q}$ y $z =\inner*{r, s}$.

	\begin{enumerate}
		\item $f\pa*{x, y} = f\pa*{y, x}$

		Se tiene lo siguiente:
		\[
			f\pa*{x, y} =\inner*{n + p, m + q} =\inner*{p + n, q + m} = f\pa*{y, x}
		\]
		\item $g\pa*{x, y} = g\pa*{y, x}$

		Se satisface lo siguiente:
		\[
			g\pa*{x, y} =\inner*{np + mq, nq + mp} =\inner*{pn + qm, pm + qn} = g\pa*{y, x}
		\]
		\item $f\pa*{x, f\pa*{y, z}} = f\pa*{f\pa*{x, y}, z}$

		Se satisface lo siguiente:
		\begin{align*}
			f\pa*{x, f\pa*{y, z}} &= f\pa*{x,\inner*{p + r, q + s}}\\
			&=\inner*{n +\pa*{p + r}, m +\pa*{q + s}}\\
			&=\inner*{\pa*{n + p} + r,\pa*{m + q} + s}\\
			&= f\pa*{\inner*{n + p, m + q}, z}\\
			&= f\pa*{f\pa*{x, y}, z}
		\end{align*}
		\item $g\pa*{x, g\pa*{y, z}} = g\pa*{g\pa*{x, y}, z}$

		Se satisface lo siguiente:
		\begin{align*}
			g\pa*{x, g\pa*{y, z}} &= g\pa*{x,\inner*{pr + qs, ps + qr}}\\
			&=\inner*{n\pa*{pr + qs} + m\pa*{ps + qr}, n\pa*{ps + qr} + m\pa*{pr + qs}}\\
			&=\inner*{npr + nqs + mps + mqr, nps + nqr + mpr + mqs}\\
			&=\inner*{\pa*{np + mq}r +\pa*{nq + mp}s,\pa*{np + mq}s +\pa*{nq + mp}r}\\
			&= g\pa*{\inner*{np + mq, nq + mp}, z}\\
			&= g\pa*{g\pa*{x, y}, z}
		\end{align*}
		\item $g\pa*{x, f\pa*{y, z}} = f\pa*{g\pa*{x, y}, g\pa*{x, z}}$

		Se satisface lo siguiente:
		\begin{align*}
			g\pa*{x, f\pa*{y, z}} &= g\pa*{x,\inner*{p + r, q + s}}\\
			&=\inner*{n\pa*{p + r} + m\pa*{q + s}, n\pa*{q + s} + m\pa*{p + r}}\\
			&=\inner*{np + nr + mq + ms, nq + ns + mp + mr}\\
			&=\inner*{\pa*{np + mq} +\pa*{nr + ms},\pa*{nq + mp} +\pa*{ns + mr}}\\
			&= f\pa*{\inner*{np + mq, nq + mp},\inner*{nr + ms, ns + mr}}\\
			&= f\pa*{g\pa*{x, y}, g\pa*{x, z}}
		\end{align*}
	\end{enumerate}
\end{proof}

\begin{lemma}
	\label{lemma: 0_enteros}
	Para cada $n\in\NN$ se satisface que $\inner*{0, 0} =_{\ZZ}\inner*{n, n}$.
\end{lemma}

\begin{proof}
	Se tiene directamente puesto que $n + 0 = 0 + n$.
\end{proof}

Con esto podemos ahora definir propiamente lo que son los números enteros en $\RCA$. Aunque lo más usual en teoría de conjuntos es usar cada clase de equivalencia como cada número entero y luego el conjunto de números enteros se defina como el conjunto cociente $\NN^{2} / =_{\ZZ}$, en $\RCA$ no existen los conjuntos de conjuntos de números naturales, por lo que no podemos irnos directamente por esa definición. En cambio, de cada clase de equivalencia podríamos elegir un único elemento como representante y luego los números enteros serían tal conjunto de representantes, el cual sí puede existir ya que es un conjunto de naturales, lo cual sí está permitido. Puesto que ya sabemos que dado un conjunto no vacío podemos determinar su elemento mínimo, entonces de cada clase utilizaremos su elemento mínimo como representante.

\begin{definition}[Números enteros]
	Definimos la función $Z:\NN^{2}\to\NN$ tal que para cada $x\in\NN^{2}$, $Z\pa*{x} =\min\bra*{x}_{=_{\ZZ}}$. Se define al conjunto de \emph{números enteros} como el conjunto $\ZZ\subseteq\NN$ tal que:
	\[
		\ZZ\coloneqq\RANGO{Z}
	\]
	es decir, un \emph{número entero} es cualquier natural $x\in\NN$ tal que existe $n\in\NN^{2}$ para el cual $x =\min\bra*{n}_{=_{\ZZ}}$.
\end{definition}

Notemos que $\ZZ$ existe por $\COMPRENSION{\Delta_{1}^{0}}$. Además, $Z$ está bien definida porque para cada $x\in\NN^{2}$ se satisface que $x\in\bra*{x}_{=_{\ZZ}}$ y luego su mínimo existe. Otra observación es que claramente para cada $x, y\in\NN^{2}$ tales que $x =_{\ZZ} y$ se satisface que $Z\pa*{x} = Z\pa*{y}$, pues $x$ y $y$ forman parte de la misma clase y luego tienen el mismo mínimo.

Ahora podemos también definir adecuadamente la suma y producto de $\ZZ$, con lo que podemos demostrar que $\ZZ$ es un anillo con unidad.

\begin{definition}[Suma y producto en $\ZZ$]
	Se define la \emph{suma} y \emph{producto} en $\ZZ$ como las siguientes funciones:
	\[
		\begin{NiceArray}{cc}
			+_{\ZZ}:\ZZ^{2}\to\ZZ &\cdot_{\ZZ}:\ZZ^{2}\to\ZZ\\
			\pa*{x, y}\mapsto Z\pa*{f\pa*{x, y}} &\pa*{x, y}\mapsto Z\pa*{g\pa*{x, y}}
		\end{NiceArray}
	\]
	donde $f$ y $g$ son las funciones de la definición \ref{def: suma_producto_z_aux}.
\end{definition}

\begin{theorem}
	\label{theo: enteros_anillo}
	El conjunto de enteros $\ZZ$ es un anillo conmutativo con unidad con las operaciones $+_{\ZZ}$ y $\cdot_{\ZZ}$.
\end{theorem}

\begin{proof}
	Sean $x, y, z\in\ZZ$, entonces mostremos cada una de las siguientes propiedades con ayuda de los lemas \ref{lemma: independencia_representantes_z}, \ref{lemma: prop_suma_prod_z} y \ref{lemma: 0_enteros}:

	\begin{enumerate}
		\item Conmutatividad para la suma
		\[
			x +_{\ZZ} y = Z\pa*{f\pa*{x, y}} = Z\pa*{f\pa*{y, x}} = y +_{\ZZ} x
		\]
		\item Asociatividad para la suma
		\begin{align*}
			x +_{\ZZ}\pa*{y +_{\ZZ} z} &= Z\pa*{f\pa*{x, Z\pa*{f\pa*{y, z}}}}\\
			&= Z\pa*{f\pa*{x, f\pa*{y, z}}}\\
			&= Z\pa*{f\pa*{f\pa*{x, y}, z}}\\
			&= Z\pa*{f\pa*{Z\pa*{f\pa*{x, y}}, z}}\\
			&=\pa*{x +_{\ZZ} y} +_{\ZZ} z
		\end{align*}
		\item Existencia del neutro aditivo

		Definamos al entero $0_{\ZZ}\coloneqq Z\pa*{\inner*{0, 0}}$, entonces si $x =\inner*{n, m}$ para algunos $n, m\in\NN$, se tiene que:
		\begin{align*}
			x +_{\ZZ} 0_{\ZZ} &= Z\pa*{f\pa*{x, 0_{\ZZ}}}\\
			&= Z\pa*{f\pa*{\inner*{n, m},Z\pa*{\inner*{0, 0}}}}\\
			&= Z\pa*{f\pa*{\inner*{n, m},\inner*{0, 0}}}\\
			&= Z\pa*{\inner*{n, m}}\\
			&= Z\pa*{x} = x
		\end{align*}
		\item Existencia de inversos aditivos

		Si $x =\inner*{n, m}$ para algunos $n, m\in\NN$, entonces definamos $-x\coloneqq Z\pa*{\inner*{m, n}}$. Luego:
		\begin{align*}
			x +_{\ZZ}\pa*{-x} &= Z\pa*{f\pa*{x, -x}}\\
			&= Z\pa*{f\pa*{\inner*{n, m}, Z\pa*{\inner*{m, n}}}}\\
			&= Z\pa*{f\pa*{\inner*{n, m},\inner*{m, n}}}\\
			&= Z\pa*{\inner*{n + m, n + m}}\\
			&= Z\pa*{\inner*{0, 0}} = 0_{\ZZ}
		\end{align*}
		\item Conmutatividad para el producto
		\[
			x\cdot_{\ZZ} y = Z\pa*{g\pa*{x, y}} = Z\pa*{g\pa*{y, x}} = y\cdot_{\ZZ} x
		\]
		\item Asociatividad para el producto
		\begin{align*}
			x\cdot_{\ZZ}\pa*{y\cdot_{\ZZ} z} &= Z\pa*{g\pa*{x, Z\pa*{g\pa*{y, z}}}}\\
			&= Z\pa*{g\pa*{x, g\pa*{y, z}}}\\
			&= Z\pa*{g\pa*{g\pa*{x, y}, z}}\\
			&= Z\pa*{g\pa*{Z\pa*{g\pa*{x, y}}, z}}\\
			&=\pa*{x\cdot_{\ZZ} y}\cdot_{\ZZ} z
		\end{align*}
		\item Existencia del neutro multiplicativo

		Definamos al entero $1_{\ZZ}\coloneqq Z\pa*{\inner*{1, 0}}$, entonces si $x =\inner*{n, m}$ para algunos $n, m\in\NN$, se tiene que:
		\begin{align*}
			x\cdot_{\ZZ} 1_{\ZZ} &= Z\pa*{g\pa*{x, 1_{\ZZ}}}\\
			&= Z\pa*{g\pa*{\inner*{n, m},Z\pa*{\inner*{1, 0}}}}\\
			&= Z\pa*{g\pa*{\inner*{n, m},\inner*{1, 0}}}\\
			&= Z\pa*{\inner*{n\cdot 1 + m\cdot 0, n\cdot 0 + m\cdot 1}}\\
			&= Z\pa*{\inner*{n, m}}\\
			&= Z\pa*{x} = x
		\end{align*}

		\item Distributividad

		\begin{align*}
			x\cdot_{\ZZ}\pa*{y +_{\ZZ} z} &= Z\pa*{g\pa*{x, Z\pa*{f\pa*{y, z}}}}\\
			&= Z\pa*{g\pa*{x, f\pa*{y, z}}}\\
			&= Z\pa*{f\pa*{g\pa*{x, y}, g\pa*{x, z}}}\\
			&= Z\pa*{f\pa*{Z\pa*{g\pa*{x, y}}, Z\pa*{g\pa*{x, z}}}}\\
			&=\pa*{x\cdot_{\ZZ} y} +_{\ZZ}\pa*{x\cdot_{\ZZ} z}
		\end{align*}
	\end{enumerate}

	Por lo tanto, $\ZZ$ es un anillo conmutativo con unidad con las operaciones $+_{\ZZ}$ y $\cdot_{\ZZ}$.
\end{proof}

Ahora, para definir a los números racionales es fundamental tener el concepto de orden en $\ZZ$. Fuera de $\RCA$, si tenemos dos números enteros $x = n - m$ y $y = p - q$ para algunos naturales $n$, $m$, $p$ y $q$, sabemos que $x < y$ si y solo si $n - m < p - q$, o equivalentemente, $n + q < p + m$. De esta forma podemos también extender el concepto de orden a $\ZZ$ en $\RCA$ como tal relación.

\begin{definition}[Orden en $\ZZ$]
	Se define la \emph{relación de orden} en $\ZZ$ como la relación $<_{\ZZ}\subseteq\ZZ^{2}$ tal que para cada $x =\inner*{n, m}, y =\inner*{p, q}\in\ZZ$:
	\[
		x <_{\ZZ} y\iff n + q < p + m
	\]
\end{definition}

Se tiene que $<_{\ZZ}$ existe por $\COMPRENSION{\Delta_{1}^{0}}$, y además es un orden lineal.

\begin{theorem}
	La relación de orden $<_{\ZZ}$ es un orden lineal en $\ZZ$.
\end{theorem}

\begin{proof}
	Sean $x =\inner*{n, m}, y =\inner*{p, q}, z =\inner*{r, s}\in\ZZ$, entonces mostremos que $<_{\ZZ}$ es irreflexiva, transitiva y tricotómica.

	\begin{enumerate}
		\item Irreflexividad

		Mostremos que $\neg\pa*{x <_{\ZZ} x}$. Puesto que $n + m = m + n$, entonces se satisface que $\neg\pa*{n + m < m + n}$, por lo que $\neg\pa*{x <_{\ZZ} x}$.

		\item Transitividad

		Supongamos que $x <_{\ZZ} y$ y $y <_{\ZZ} z$, entonces $n + q < p + m$ y $p + s < r + q$, por lo que se tiene lo siguiente:
		\begin{align*}
			\pa*{n + q} +\pa*{p + s} &<\pa*{p + m} +\pa*{r + q}\\
			\pa*{n + s} +\pa*{p + q} &<\pa*{r + m} +\pa*{p + q}\\
			n + s &< r + m
		\end{align*}
		Por lo tanto, $x <_{\ZZ} z$.

		\item Tricotomía

		Como $n, m, p, q\in\NN$, y el orden en $\NN$ es tricotómico, entonces se satisface una y solo una de las siguientes afirmaciones: $n + q = p + m$, $n + q < p + m$ o $n + q > p + m$.

		Si $n + q = p + m$, entonces $x =_{\ZZ} y$, lo cual implicaría que $x = y$ ya que al ser enteros, serían representantes de la misma clase de equivalencia.

		Si $n + q < p + m$, entonces $x <_{\ZZ} y$ por definición, y si $n + q > p + m$, también por definición $y <_{\ZZ} x$.

		Por lo tanto, $<_{\ZZ}$ es tricotómica.
	\end{enumerate}

	Por lo tanto, $<_{\ZZ}$ es un orden lineal en $\ZZ$.
\end{proof}

Para finalizar la construcción de $\ZZ$, mostremos que los números naturales se pueden encajar en $\ZZ$, lo cual podemos entender como que existe un subconjunto de $\ZZ$ que se comporta de manera idéntica como lo haría $\NN$ en cuanto a la suma, producto y orden. Formalmente eso significa que debe existir una función inyectiva de $\NN$ a $\ZZ$ que conserve la suma, producto y orden de $\NN$ a $\ZZ$.

Puesto que nosotros ya interpretábamos cada par $\inner*{n, m}$ de naturales como la resta $n - m$, entonces los números naturales serían los enteros tales que su segunda componente es $0$. Por ejemplo, el número $2$ en $\ZZ$ sería el representante de la clase de $\inner*{2, 0}$, pues $2 = 2 - 0$.

Aseguramos que esta traducción de naturales a enteros es aquella función inyectiva que nos permite encajar los naturales en los enteros.

\begin{theorem}[Encaje de $\NN$ en $\ZZ$]
	Defínase la función $N:\NN\to\ZZ$ tal que para cada $n\in\NN$, $N\pa*{n}\coloneqq Z\pa*{\inner*{n, 0}}$. Entonces $N$ es inyectiva y satisface las siguientes propiedades:

	\begin{enumerate}
		\item Para cada $n, m\in\NN$, $N\pa*{n + m} = N\pa*{n} +_{\ZZ} N\pa*{m}$.
		\item Para cada $n, m\in\NN$, $N\pa*{nm} = N\pa*{n}\cdot_{\ZZ} N\pa*{m}$.
		\item Para cada $n, m\in\NN$, $n < m$ si y solo si $N\pa*{n} <_{\ZZ} N\pa*{m}$.
	\end{enumerate}
\end{theorem}

\begin{proof}
	Sean $n, m\in\NN$, entonces:

	\begin{enumerate}
		\item Inyectividad de $N$

		Supongamos que $N\pa*{n} = N\pa*{m}$, entonces se tiene que $Z\pa*{\inner*{n, 0}} = Z\pa*{\inner*{m, 0}}$, lo cual significa que $\inner*{n, 0} =_{\ZZ}\inner*{m, 0}$, es decir que $n + 0 = 0 + m$, o equivalentemente, $n = m$. Por lo tanto, $N$ es inyectiva.

		\item Conservación de la suma

		Se tiene lo siguiente:
		\begin{align*}
			N\pa*{n + m} &= Z\pa*{\inner*{n + m, 0}}\\
			&= Z\pa*{f\pa*{\inner*{n, 0},\inner*{m, 0}}}\\
			&= Z\pa*{f\pa*{Z\pa*{\inner*{n, 0}}, Z\pa*{\inner*{m, 0}}}}\\
			&= Z\pa*{f\pa*{N\pa*{n}, N\pa*{m}}}\\
			&= N\pa*{n} +_{\ZZ} N\pa*{m}
		\end{align*}

		\item Conservación del producto

		Se tiene lo siguiente:
		\begin{align*}
			N\pa*{nm} &= Z\pa*{\inner*{nm, 0}}\\
			&= Z\pa*{\inner*{nm + 0\cdot 0, n\cdot 0 + 0\cdot m}}\\
			&= Z\pa*{g\pa*{\inner*{n, 0},\inner*{m, 0}}}\\
			&= Z\pa*{g\pa*{Z\pa*{\inner*{n, 0}}, Z\pa*{\inner*{m, 0}}}}\\
			&= Z\pa*{g\pa*{N\pa*{n}, N\pa*{m}}}\\
			&= N\pa*{n}\cdot_{\ZZ} N\pa*{m}
		\end{align*}

		\item Conservación del orden

		Sean $N\pa*{n} =\inner*{p, q}$ y $N\pa*{m} =\inner*{r, s}$ para algunos $p, q, r, s\in\NN$, entonces sabemos que $n + q = p$ y $m + s = r$ por definición de $N$, entonces:
		\begin{align*}
			n < m &\iff n + q < m + q\\
			&\iff p < m + q\\
			&\iff p + s < m + q + s\\
			&\iff p + s < r + q\\
			&\iff\inner*{p, q} <_{\ZZ}\inner*{r, s}\\
			&\iff N\pa*{n} <_{\ZZ} N\pa*{m}
		\end{align*}
	\end{enumerate}
\end{proof}

Con esto hemos mostrado que al trabajar con los enteros, también podemos hacer referencia a los naturales si es que necesitamos tratar con ellos. Aunque cuando hablamos de los naturales en $\ZZ$ no necesariamente estamos realizando las mismas operaciones en el propio $\NN$, es decir, no necesariamente $2_{\ZZ} +_{\ZZ} 3_{\ZZ} = 2 + 3$, la inyección $N$ nos permite traducir cada expresión de un conjunto a otro de manera sencilla, indicándonos que podemos ignorar esa diferencia entre $\NN$ y $\RANGO{N}$, y en cambio tratarlos como el mismo conjunto. Por eso mismo, no haremos distinción simbólica entre $+_{\ZZ}$ y la suma en $\NN$ al trabajar en $\ZZ$ debido a ese encaje.

\subsection{Los números racionales}

Ya que hemos construido los números enteros, ya directamente podemos definir lo que son los números racionales, lo cuales son los más fundamentales para el estudio del Análisis Real que realizaremos.

Fuera de $\RCA$, los números racionales se caracterizan por tener \emph{división}, y en general cualquier número racional $q$ puede expresarse como un cociente entre números enteros $q = n / m$, donde $m > 0$. De manera similar podemos codificar un número racional como un par de números enteros interpretados de la forma anterior. Por ejemplo, el número $2/3$ sería el par $(2, 3)$. Sin embargo, el par no es único, pues $2/3$ también puede representarse con los pares $(4, 6)$, $(6, 9)$, $(8, 12)$, etc. Entonces todos estos pares se deberían considerar \emph{equivalentes} al interpretarse como el mismo racional. En general, para cada par de pares $(p, q)$ y $(r, s)$ que representan a los racionales $p/q$ y $r/s$, estos son equivalentes si y solo si $p/q = r/s$, o equivalentemente $ps = rq$. De esta manera podemos definir la relación de equivalencia entre pares de enteros en $\RCA$ que codifica la igualdad entre racionales.

\begin{definition}[Igualdad en $\QQ$]
	Se define la relación $=_{\QQ}\subseteq\pa*{\ZZ\times\ZZ^{+}}\times\pa*{\ZZ\times\ZZ^{+}}$ como sigue para cada $n, m, p, q\in\NN$:
	\[
		\inner*{n, m} =_{\QQ}\inner*{p, q}\iff nq = pm
	\]
\end{definition}

Aquí $\ZZ^{+}$ se entiende como el conjunto de enteros positivos, es decir el conjunto $\ZZ^{+} =\set*{n:\pa*{n\in\ZZ}\land\pa*{n >_{\ZZ} 0}}$. Se tiene también que esta relación existe por $\COMPRENSION{\Delta_{1}^{0}}$.

De manera análoga al teorema \ref{theo: relacion_equivalencia_z} se demuestra que $=_{\QQ}$ es de equivalencia, por lo que para cada $q\in\ZZ\times\ZZ^{+}$ existe su clase de equivalencia $\bra*{q}_{=_{\QQ}}$ por $\COMPRENSION{\Delta_{1}^{0}}$. Entonces, de manera similar definiremos a los números racionales como el conjunto de elementos mínimos de cada clase de equivalencia.

\begin{definition}[Números racionales]
	Defínase la función $Q:\ZZ\times\ZZ^{+}\to\NN$ tal que para cada $q\in\ZZ\times\ZZ^{+}$, $Q\pa*{q} =\min\bra*{q}_{=_{\QQ}}$. Entonces definimos al \emph{conjunto de números racionales} como el subconjunto $\QQ\subseteq\NN$ tal que:
	\[
		\QQ\coloneqq\RANGO{Q}
	\]
	es decir, un \emph{número racional} es cualquier natural $q\in\NN$ tal que existe $n\in\ZZ\times\ZZ^{+}$ para el cual $q =\min\bra*{n}_{=_{\QQ}}$.
\end{definition}

De manera análoga podemos definir las operaciones en $\QQ$ por medio de dos funciones auxiliares que codifican la suma y producto de racionales usando los pares de enteros. Fuera de $\RCA$, si tenemos dos racionales $x = n/m$ y $y = p/q$, entonces obtenemos lo siguiente:
\[
	\begin{NiceArray}{c}
		x + y =\dfrac{n}{m} +\dfrac{p}{q} =\dfrac{nq + mp}{mq}\\
		xy =\pa*{\dfrac{n}{m}}\pa*{\dfrac{p}{q}} =\dfrac{np}{mq}
	\end{NiceArray}
\]

Por lo que la suma y producto estarían codificados como los pares $\pa*{nq + mp, mq}$ y $\pa*{np, mq}$ respectivamente, lo cual podemos llevar a $\RCA$ de la siguiente forma.

\begin{definition}
	\label{def: suma_producto_q_aux}
	Definimos las siguientes funciones:
	\[
		\begin{NiceArray}{cc}
			f:\pa*{\ZZ\times\ZZ^{+}}\times\pa*{\ZZ\times\ZZ^{+}}\to\NN & g:\pa*{\ZZ\times\ZZ^{+}}\times\pa*{\ZZ\times\ZZ^{+}}\to\NN\\
			\pa*{\inner*{n, m},\inner*{p, q}}\mapsto\inner*{nq + mp, mq} &\pa*{\inner*{n, m},\inner*{p, q}}\mapsto\inner*{np, mq}
		\end{NiceArray}
	\]
\end{definition}

Al igual que en caso de los enteros, es necesario asegurar que estas operaciones son independientes del representante bajo la igualdad en $\QQ$.

\begin{lemma}
	\label{lemma: independencia_representantes_q}
	Para cada $x, y, z, w\in\ZZ\times\ZZ^{+}$ se satisface que si $x =_{\QQ} z$ y $y =_{\QQ} w$, entonces $f\pa*{x, y} =_{\QQ} f\pa*{z, w}$ y $g\pa*{x, y} =_{\QQ} g\pa*{z, w}$.
\end{lemma}

\begin{proof}
	Sean $n, p, r, t\in\ZZ$ y $m, q, s, u\in\ZZ^{+}$ tales que $x =\inner*{n, m}$, $y =\inner*{p, q}$, $z =\inner*{r, s}$ y $w =\inner*{t, u}$ donde además $x =_{\QQ} z$ y $y =_{\QQ} w$. Se tiene entonces que $ns = rm$ y $pu = tq$.

	\begin{enumerate}
		\item Para $f$ se tiene lo siguiente:
		\begin{align*}
			\pa*{nq + mp}\pa*{su} &= nqsu + mpsu\\
			&= rmqu + tqms\\
			&=\pa*{ru + st}\pa*{mq}
		\end{align*}
		Por lo que $f\pa*{x, y} =_{\QQ} f\pa*{z, w}$.
		\item Para $g$ se tiene lo siguiente:
		\[
			\pa*{np}\pa*{su} =\pa*{ns}\pa*{pu} =\pa*{rm}\pa*{tq}
		\]
		Por lo que $g\pa*{x, y} =_{\QQ} g\pa*{z, w}$.
	\end{enumerate}
\end{proof}

Una observación a realizar es que en los racionales se cumple la cancelación en la división, es decir, si nosotros tenemos el racional $np/nq$ donde $n\neq 0$, entonces $np/nq = p/q$. Traducido a la construcción que estamos realizando se obtiene el siguiente teorema.

\begin{lemma}
	\label{lemma: cancelacion_cociente}
	Para cada $p\in\ZZ$ y cada $n, q\in\ZZ^{+}$ se satisface que $\inner*{np, nq} =_{\QQ}\inner*{p, q}$.
\end{lemma}

\begin{proof}
	Se sigue directamente puesto que $\pa*{np}q = p\pa*{nq}$.
\end{proof}

Con esto podemos demostrar que estas funciones se comportan bien, lo cual significa que son conmutativas, asociativas y distributivas.

\begin{lemma}
	\label{lemma: prop_suma_prod_q}
	Para cada $x, y, z\in\ZZ\times\ZZ^{+}$ se satisface lo siguiente:

	\begin{enumerate}
		\item $f\pa*{x, y} = f\pa*{y, x}$
		\item $g\pa*{x, y} = g\pa*{y, x}$
		\item $f\pa*{x, f\pa*{y, z}} = f\pa*{f\pa*{x, y}, z}$
		\item $g\pa*{x, g\pa*{y, z}} = g\pa*{g\pa*{x, y}, z}$
		\item $g\pa*{x, f\pa*{y, z}} =_{\QQ} f\pa*{g\pa*{x, y}, g\pa*{x, z}}$
	\end{enumerate}
\end{lemma}

\begin{proof}
	Las demostraciones de los primeros cuatro puntos son análogas a lo que se realiza en las demostraciones del teorema \ref{lemma: prop_suma_prod_z}. La única necesaria a demostrar es la del último punto que es la distributividad.

	Sean $n, p, r\in\ZZ$ y $m, q, s\in\ZZ^{+}$ tales que $x =\inner*{n, m}$, $y =\inner*{p, q}$, $z =\inner*{r, s}$, entonces obtenemos lo siguiente usando el lema \ref{lemma: cancelacion_cociente}:
	\begin{align*}
		g\pa*{x, f\pa*{y, z}} &= g\pa*{x,\inner*{ps + qr, qs}}\\
		&=\inner*{n\pa*{ps + qr}, m\pa*{qs}}\\
		&=\inner*{nps + nqr, mqs}\\
		&=_{\QQ}\inner*{m\pa*{nps + nqr}, m^{2}qs}\\
		&=\inner*{\pa*{np}\pa*{ms} +\pa*{mq}\pa*{nr},\pa*{mq}\pa*{ms}}\\
		&= f\pa*{\inner*{np, mq},\inner*{nr, ms}}\\
		&= f\pa*{g\pa*{x, y}, g\pa*{x, z}}
	\end{align*}
\end{proof}

\begin{lemma}
	\label{lemma: 0_1_racionales}
	Para cada $n\in\ZZ^{+}$ se satisface que $\inner*{0, 1} =_{\QQ}\inner*{0, n}$ y $\inner*{1, 1} =_{\QQ}\inner*{n, n}$.
\end{lemma}

\begin{proof}
	Se sigue directamente de las igualdades $0n = 0$ y $n = n$.
\end{proof}

Podemos entonces ahora definir las operaciones de suma y producto en $\QQ$ de manera coherente.

\begin{definition}[Suma y producto en $\QQ$]
	Se define la \emph{suma} y \emph{producto} en $\QQ$ como las siguientes funciones:
	\[
		\begin{NiceArray}{cc}
			+_{\QQ}:QQ^{2}\to\QQ &\cdot_{\QQ}:QQ^{2}\to\QQ\\
			\pa*{x, y}\mapsto Q\pa*{f\pa*{x, y}} &\pa*{x, y}\mapsto Q\pa*{g\pa*{x, y}}
		\end{NiceArray}
	\]
	donde $f$ y $g$ son las funciones de la definición \ref{def: suma_producto_q_aux}.
\end{definition}

Se puede entonces probar una propiedad fundamental la cual es que $\QQ$ es un campo.

\begin{theorem}
	El conjunto de números racionales $\QQ$ es un campo con las operaciones $+_{\QQ}$ y $\cdot_{\QQ}$.
\end{theorem}

\begin{proof}
	Solo es necesario demostrar la existencia de neutros e inversos (aditivos y multiplicativos) en $\QQ$, pues las otras propiedades de campo son análogas a las de la demostración del teorema \ref{theo: enteros_anillo}.

	\begin{enumerate}
		\item Neutro aditivo

		Sea $x =\inner*{p, q}\in\QQ$ para algunos $p\in\ZZ$ y $q\in\ZZ^{+}$. Entonces si definimos $0_{\QQ}\coloneqq Q\pa*{\inner*{0, 1}}$, se obtiene lo siguiente:
		\begin{align*}
			x +_{\QQ} 0_{\QQ} &= Q\pa*{f\pa*{x, 0_{\QQ}}}\\
			&= Q\pa*{f\pa*{\inner*{p, q}, Q\pa*{\inner*{0, 1}}}}\\
			&= Q\pa*{f\pa*{p, q},\inner*{0, 1}}\\
			&= Q\pa*{p\cdot 1 + q\cdot 0, q\cdot 1}\\
			&= Q\pa*{\inner*{p, q}}\\
			&= Q\pa*{x} = x
		\end{align*}
		\item Neutro multiplicativo

		Sea $x =\inner*{p, q}\in\QQ$ para algunos $p\in\ZZ$ y $q\in\ZZ^{+}$. Entonces si definimos $1_{\QQ}\coloneqq Q\pa*{\inner*{1, 1}}$, se obtiene lo siguiente:
		\begin{align*}
			x\cdot_{\QQ} 1_{\QQ} &= Q\pa*{g\pa*{x, 1_{\QQ}}}\\
			&= Q\pa*{g\pa*{\inner*{p, q}, Q\pa*{\inner*{1, 1}}}}\\
			&= Q\pa*{g\pa*{p, q},\inner*{1, 1}}\\
			&= Q\pa*{\inner*{p, q}}\\
			&= Q\pa*{x} = x
		\end{align*}

		\item Inverso aditivo

		Sea $x =\inner*{p, q}\in\QQ$ para algunos $p\in\ZZ$ y $q\in\ZZ^{+}$. Entonces si definimos $-x\coloneqq Q\pa*{\inner*{-p, q}}$, se obtiene lo siguiente:
		\begin{align*}
			x +_{\QQ}\pa*{-x} &= Q\pa*{f\pa*{x, -x}}\\
			&= Q\pa*{f\pa*{\inner*{p, q}, Q\pa*{\inner*{-p, q}}}}\\
			&= Q\pa*{f\pa*{p, q},\inner*{-p, q}}\\
			&= Q\pa*{pq + q\pa*{-p}, q^{2}}\\
			&= Q\pa*{\inner*{0, q^{2}}}\\
			&= Q\pa*{\inner*{0, 1}} = 0_{\QQ}
		\end{align*}

		\item Inverso multiplicativo

		Sea $x =\inner*{p, q}\in\QQ$ para algunos $p\in\ZZ$ y $q\in\ZZ^{+}$. Si $x\neq 0_{\QQ}$, entonces $p\neq 0$, luego tenemos dos casos: $p > 0$ o $p < 0$.

		Si $p > 0$, entonces definimos $x^{-1}\coloneqq Q\pa*{\inner*{q, p}}$. Luego:
		\begin{align*}
			x\cdot_{\QQ} x^{-1} &= Q\pa*{g\pa*{x, x^{-1}}}\\
			&= Q\pa*{g\pa*{\inner*{p, q}, Q\pa*{\inner*{q, p}}}}\\
			&= Q\pa*{g\pa*{p, q},\inner*{q, p}}\\
			&= Q\pa*{\inner*{pq, pq}}\\
			&= Q\pa*{\inner*{1, 1}} = 1_{\QQ}
		\end{align*}

		En cambio, si $p < 0$, entonces definimos $x^{-1}\coloneqq Q\pa*{\inner*{-q, -p}}$. Luego:
		\begin{align*}
			x\cdot_{\QQ} x^{-1} &= Q\pa*{g\pa*{x, x^{-1}}}\\
			&= Q\pa*{g\pa*{\inner*{p, q}, Q\pa*{\inner*{-q, -p}}}}\\
			&= Q\pa*{g\pa*{p, q},\inner*{-q, -p}}\\
			&= Q\pa*{\inner*{-pq, -pq}}\\
			&= Q\pa*{\inner*{1, 1}} = 1_{\QQ}
		\end{align*}

		En cualquier caso existe $x^{-1}\in\QQ$ tal que $x\cdot_{\QQ}x^{-1} = 1_{\QQ}$
	\end{enumerate}

	Por lo tanto, el conjunto $\QQ$ es un campo con las operaciones $+_{\QQ}$ y $\cdot_{\QQ}$.
\end{proof}

Para terminar con la construcción de los racionales definamos el orden en ellos. Sean $x = n / m$ y $y = p / q$. Puesto que $m$ y $p$ se toman como enteros positivos, entonces se tiene que $x < y$ si y solo si $n / m < p / q$, o equivalentemente $nq < pm$. Entonces de esta manera podemos ahora caracterizar el orden de $\QQ$ en $\RCA$.

\begin{definition}[Orden en $\QQ$]
	Se define la \emph{relación de orden} en $\QQ$ como la relación binaria $<_{\QQ}\subseteq\QQ^{2}$ tal que para cada $x =\inner*{n, m}, y =\inner*{p, q}\in\QQ$ para algunos $n, m, p, q\in\ZZ$:
	\[
		x <_{\QQ} y\iff nq < pm
	\]
\end{definition}

\begin{theorem}
	La relación $<_{\QQ}$ es un orden lineal.
\end{theorem}

\begin{proof}
	Sean $x =\inner*{n, m}, y =\inner*{p, q}, z =\inner*{r, s}\in\QQ$ para algunos $n, p, r\in\ZZ$ y $m, q, s\in\ZZ^{+}$. Mostremos entonces que $<_{\QQ}$ es irreflexiva, transitiva y tricotómica.

	\begin{enumerate}
		\item Irreflexividad

		Puesto que el orden en $\ZZ$ es irreflexiva, sabemos entonces que $\neg\pa*{nm < nm}$, es decir, $\neg\pa*{x <_{\QQ} x}$.

		\item Transitividad

		Supongamos que $x <_{\QQ} y$ y $y <_{\QQ} z$, entonces $nq < pm$ y $ps < rq$. Puesto que $s > 0$, entonces $nqs < pms$. Como $m > 0$, entonces $psm < rqm$. Por transitividad del orden en $\ZZ$ tenemos que $nqs < rqm$, y como $q > 0$, entonces $ns < rm$. Por lo tanto, $x <_{\QQ} z$.

		\item Tricotomía

		Por la tricotomía en $\ZZ$ sabemos que se satisface una y solo una de las siguientes afirmaciones: $nq = pm$, $nq < pm$ o $nq > pm$.

		Si $nq = pm$, entonces $x =_{\QQ} y$, y puesto que $x, y\in\QQ$, entonces $x = y$.

		Si $nq < pm$, entonces $x <_{\QQ} y$; y si $nq > pm$, entonces $y <_{\QQ} x$.
	\end{enumerate}

	Por lo tanto, $<_{\QQ}$ es un orden lineal.
\end{proof}

Finalmente podemos demostrar igualmente que podemos encajar $\ZZ$ en $\QQ$ mediante el encaje $n\mapsto Q\pa*{\inner*{n, 1}}$, pues se interpreta que $n = n / 1$.

\begin{theorem}[Encaje de $\ZZ$ en $\QQ$]
	Defínase la función $Z:\ZZ\to\QQ$ tal que para cada $n\in\ZZ$, $Z\pa*{n}\coloneqq Q\pa*{\inner*{n, 1}}$. Entonces $Z$ es inyectiva y satisface las siguientes propiedades:

	\begin{enumerate}
		\item Para cada $n, m\in\ZZ$, $Z\pa*{n + m} = Z\pa*{n} +_{\QQ} Z\pa*{m}$.
		\item Para cada $n, m\in\ZZ$, $Z\pa*{nm} = Z\pa*{n}\cdot_{\QQ} Z\pa*{m}$.
		\item Para cada $n, m\in\ZZ$, $n < m$ si y solo si $Z\pa*{n} <_{\QQ} Z\pa*{m}$.
	\end{enumerate}
\end{theorem}

\begin{proof}
	Sean $p, q\in\ZZ$, entonces:

	\begin{enumerate}
		\item Inyectividad de $Z$

		Supongamos que $Z\pa*{p} = Z\pa*{q}$, entonces se tiene que $Q\pa*{\inner*{p, 1}} = Q\pa*{\inner*{q, 1}}$, lo cual significa que $\inner*{p, 1} =_{\QQ}\inner*{q, 1}$, es decir que $p\cdot 1 = q\cdot 1$, o equivalentemente, $p = q$. Por lo tanto, $Z$ es inyectiva.

		\item Conservación de la suma

		Se tiene lo siguiente:
		\begin{align*}
			Z\pa*{p + q} &= Q\pa*{\inner*{p + q, 1}}\\
			&= Q\pa*{f\pa*{\inner*{p, 1},\inner*{q, 1}}}\\
			&= Q\pa*{f\pa*{Q\pa*{\inner*{p, 1}}, Q\pa*{\inner*{q, 1}}}}\\
			&= Q\pa*{f\pa*{Z\pa*{n}, Z\pa*{m}}}\\
			&= Z\pa*{n} +_{\QQ} Z\pa*{m}
		\end{align*}

		\item Conservación del producto

		Se tiene lo siguiente:
		\begin{align*}
			Z\pa*{pq} &= Q\pa*{\inner*{pq, 1}}\\
			&= Q\pa*{g\pa*{\inner*{p, 1},\inner*{q, 1}}}\\
			&= Q\pa*{g\pa*{Q\pa*{\inner*{p, 1}}, Q\pa*{\inner*{q, 1}}}}\\
			&= Q\pa*{g\pa*{Z\pa*{p}, Z\pa*{q}}}\\
			&= Z\pa*{n}\cdot_{\QQ} Z\pa*{m}
		\end{align*}

		\item Conservación del orden

		Sean $Z\pa*{p} =\inner*{n, m}$ y $Z\pa*{q} =\inner*{r, s}$ para algunos $n, r\in\ZZ$ y $m, s\in\ZZ^{+}$, entonces sabemos que $pm = n$ y $qs = r$ por definición de $Z$, entonces:
		\begin{align*}
			p < q &\iff pm < qm\\
			&\iff n < mq\\
			&\iff ns < mqs\\
			&\iff ns < mr\\
			&\iff\inner*{n, m} <_{\QQ}\inner*{r, s}\\
			&\iff Z\pa*{p} <_{\QQ} Z\pa*{q}
		\end{align*}
	\end{enumerate}
\end{proof}

\subsection{Los números reales}

Con los racionales a nuestra disposición, el siguiente paso es definir qué significa ser un número real. Aquí aparece una primera dificultad: los reales no pueden formar un conjunto de $\RCA$, pues cualquier construcción de $\RR$ produce un conjunto incontable y, por lo tanto, imposible de codificar como subconjunto de $\NN$. Aun cuando existan modelos incontables de $\RCA$, su estructura interna es \emph{contable}, así que no hay manera de representar a $\RR$ como un conjunto dentro del sistema. Lo que haremos, entonces, será definir qué es \emph{un} número real: cada real individual sí podrá codificarse, aunque la totalidad de ellos no.

La construcción que seguiremos es la de Cauchy. Para motivarla, recordemos (fuera de $\RCA$) que $\QQ$, con la métrica inducida por el valor absoluto, no es un espacio métrico completo: existen sucesiones de Cauchy de racionales que no convergen en $\QQ$. Por ejemplo, la sucesión $\set*{x_{n}}_{n\in\omega}$ definida por $x_{0} = 1$ y $x_{n + 1} = x_{n} / 2 + 1 / x_{n}$ es, a partir de su segundo término, decreciente y acotada inferiormente por $\sqrt{2}$, y por lo tanto convergente en $\RR$. Como toda sucesión convergente en $\RR$ es de Cauchy, se trata de una sucesión de Cauchy de números racionales; sin embargo, $\lim_{n\to\infty}x_{n} =\sqrt{2}\notin\QQ$.

La idea es entonces construir $\RR$ \emph{completando} a $\QQ$: consideramos todas las sucesiones de Cauchy de racionales e identificamos aquellas que deberían tener el mismo límite. Como ese límite no necesariamente existe en $\QQ$, no podemos compararlo directamente; lo que sí podemos hacer es examinar la sucesión resta $\set*{x_{n} - y_{n}}_{n\in\omega}$, la cual debería converger a $0$. Así, diremos que dos sucesiones de Cauchy $\set*{x_{n}}_{n\in\omega},\set*{y_{n}}_{n\in\omega}\subseteq\QQ$ son \emph{equivalentes} si para cada $\varepsilon > 0$ existe $N\in\omega$ tal que $\abs*{x_{n} - y_{n}} <\varepsilon$ para cada $n\geq N$.

Esta definición, sin embargo, tiene un costo: vista como fórmula es $\Pi_{2}^{0}$, y trabajar en $\RCA$ con una noción de igualdad tan compleja resulta poco conveniente, por lo que nos interesa simplificarla. Observemos qué dice realmente la definición: la resta se ``acerca'' al cero tanto como queramos (el papel de $\varepsilon$) a partir de cierto índice (el papel de $N$), de modo que $N$ es, implícitamente, una función $N\pa*{\varepsilon}$. Una primera simplificación es no usar un $\varepsilon$ arbitrario, sino una sucesión concreta de valores que decaiga a cero; una buena propuesta son los números $2^{1 - n}$ con $n\in\omega$, que decaen rápidamente a $0$ conforme $n$ crece. Con esto, $N$ pasa a ser una función de $n$. Aun así, la función $N\pa*{n}$ puede comportarse de manera errática: puede ser mayor que $n$ para algunos valores y menor para otros, de modo que determinarla puede ser complicado y variar de una pareja de sucesiones a otra. La segunda simplificación es entonces eliminar por completo la búsqueda de $N$ imponiendo $N\pa*{n} = n$.

Llegamos así a la siguiente definición: dos sucesiones $\set*{x_{n}}_{n\in\omega},\set*{y_{n}}_{n\in\omega}\subseteq\QQ$ son equivalentes si para cada $n\in\omega$ se satisface que $\abs*{x_{n} - y_{n}} < 2^{1 - n}$. Ahora bien, esta condición obliga a que la sucesión resta decaiga a un ritmo exponencial, y no toda sucesión de Cauchy decae tan rápido. Para que la noción de equivalencia sea coherente, basta restringirnos desde el inicio a sucesiones de Cauchy que decaigan a ese mismo ritmo: pedimos que para cada $k, n, m\in\omega$, si $n, m\geq k$, entonces $\abs*{x_{n} - x_{m}} < 2^{-k}$. Esta condición se simplifica tomando $n = k$ y escribiendo $m = k + p$ con $p\in\omega$: trabajaremos entonces con las sucesiones $\set*{x_{n}}_{n\in\omega}\subseteq\QQ$ tales que para cada $k, p\in\omega$, $\abs*{x_{k} - x_{k + p}} < 2^{-k}$. Con todo esto podemos definir qué significa ser un número real en $\RCA$.

\begin{definition}[Sucesiones]
	Una \emph{sucesión de números racionales} es una función $f:\NN\to\QQ$. Se denotará a la sucesión $f$ por $\set*{x_{k}}_{k\in\NN}$ si para cada $k\in\NN$, $f\pa*{k} = x_{k}$.
\end{definition}

\begin{definition}[Números reales]
	Un \emph{número real} es cualquier sucesión de racionales $\set*{x_{k}}_{k\in\NN}$ tal que para cada $k, p\in\NN$, $\abs*{x_{k} - x_{k + p}} < 2^{-k}$.
\end{definition}

\begin{definition}[Igualdad entre números reales]
	Sean $x =\set*{x_{k}}_{k\in\NN}$ y $y =\set*{y_{k}}_{k\in\NN}$ números reales. Decimos que $x$ y $y$ son \emph{iguales}, denotado por $x =_{\RR} y$, si para cada $k$ se satisface que $\abs*{x_{k} - y_{k}} < 2^{-k + 1}$.
\end{definition}

Aunque esta relación $=_{\RR}$ no es un conjunto que exista en $\RCA$, ya que es una relación entre conjuntos y no entre números, sí podemos verificar que satisface las propiedades de una relación de equivalencia.

\begin{theorem}
	Se demuestra lo siguiente:

	\begin{enumerate}
		\item Para cada número real $x$ se satisface que $x =_{\RR} x$.
		\item Para cada par de reales $x$ y $y$ se satisface que $x =_{\RR} y$ implica $y =_{\RR} x$.
		\item Para cada terna de reales $x$, $y$ y $z$ se satisface que $x =_{\RR} y$ y $y =_{\RR} z$ implica $x =_{\RR} z$.
	\end{enumerate}
\end{theorem}

\begin{proof}
	Consideremos a los reales $x =\set*{x_{k}}_{k\in\NN}$, $y =\set*{y_{k}}_{k\in\NN}$ y $z =\set*{z_{k}}_{k\in\NN}$.

	\begin{enumerate}
		\item Se satisface claramente para cada $k$ que:
		\[
			\abs*{x_{k} - x_{k}} = 0 < 2^{-k + 1}
		\]
		Por lo que $x =_{\RR} x$.
		\item Si $x =_{\RR} y$, entonces sabemos que para cada $k$ se satisface que $\abs*{x_{k} - y_{k}} < 2^{-k + 1}$. Directamente se sigue que para cada $k$, $\abs*{y_{k} - x_{k}} < 2^{-k + 1}$, es decir, $y =_{\RR} x$.
		\item Si $x =_{\RR} y$ y $y =_{\RR} z$, entonces por definición para cada $k$, $\abs*{x_{k} - y_{k}} < 2^{-k + 1}$ y $\abs*{y_{k} - z_{k}} < 2^{-k + 1}$. Consideremos un valor fijo $p > 1$, entonces por la desigualdad del triángulo y al ser $x$, $y$ y $z$ números reales se sigue que para cada $k$:
		\begin{align*}
			\abs*{x_{k} - z_{k}} &\leq\abs*{x_{k} - x_{k + p}} +\abs*{x_{k + p} - y_{k + p}} +\abs*{y_{k} - y_{k + p}} +\abs*{y_{k} - z_{k}}\\
			&< 2^{-k} + 2^{- k - p + 1} + 2^{-k} - 2^{-k + 1}\\
			&= 2^{- k + 1}\pa*{2^{-1} + 2^{-p} + 2^{-1} + 1}\\
			&= 2^{-k + 1} 2^{-p + 1}\\
			&< 2^{-k + 1}
		\end{align*}
	\end{enumerate}
\end{proof}

Ahora definamos el orden dentro de los reales. Fuera de $\RCA$, pensemos en dos números reales $\set*{x_{n}}_{n\in\omega}$ y $\set*{y_{k}}_{k\in\omega}$. Consideremos que $x =\lim_{n\to\infty}x_{n}$ y $y =\lim_{n\to\infty}y_{n}$ son tales que $x\leq y$. Por la propuesta de definición que tenemos de convergencia en $\RCA$ sabemos que para cada $n, p\in\omega$:
\[
	\begin{NiceArray}{c}
		\abs*{x_{n} - x_{n + p}} < 2^{-n}\\
		\abs*{y_{n} - y_{n + p}} < 2^{-n}
	\end{NiceArray}
\]

Tomando el límite cuando $p\to\infty$ notemos que se tiene lo siguiente para cada $n\in\omega$:
\[
	\begin{NiceArray}{c}
		\abs*{x_{n} - x} < 2^{-n}\\
		\abs*{y_{n} - y} < 2^{-n}
	\end{NiceArray}
\]
de lo cual se sigue que $x_{n} - 2^{-n} < x$ y $y < y_{n} + 2^{n}$. Puesto que $x\leq y$ se tiene que para cada $n\in\omega$, $x_{n} - 2^{-n} < x\leq y < y_{n} + 2^{-n}$, o equivalentemente se satisface que $x_{n} < y_{n} + 2^{-n + 1}$. Además, notemos que tomando el límite cuando $n\to\infty$ obtenemos que:
\[
	\lim_{n\to\infty} x_{n}\leq\lim_{n\to\infty}\pa*{y_{n} + 2^{-n + 1}}\implies x\leq y
\]

Por lo tanto, podemos concluir que $x\leq y$ si y solo si para cada $n\in\omega$ se satisface que $x_{n} < y_{n} + 2^{-n + 1}$. De esta forma podemos definir la relación de orden de los reales en $\RCA$.

\begin{definition}[Orden en $\RR$]
	Sean $x =\set*{x_{k}}_{k\in\NN}$ y $y =\set*{y_{k}}_{k\in\NN}$ números reales.

	\begin{itemize}
		\item Denotaremos por $x\leq_{\RR} y$ si para cada $k$ se satisface que $x_{k} < y_{k} + 2^{-k + 1}$.

		\item Denotaremos por $x <_{\RR} y$ si no se satisface que $y\leq_{\RR}x$.
	\end{itemize}
\end{definition}

Mostremos que de esta definición de orden se puede demostrar que se satisfacen las propiedades de una relación de orden lineal.

\begin{theorem}
	Se demuestra lo siguiente:

	\begin{enumerate}
		\item Si $x$ es un real, entonces se satisface que $x\leq_{\RR} x$.
		\item Si $x$, $y$ y $z$ son números reales, $x\leq_{\RR} y$ y $y\leq_{\RR} z$, entonces $x\leq_{\RR} z$.
		\item Si $x$ y $y$ son números reales, $x\leq_{\RR} y$ y $y\leq_{\RR} x$, entonces $x =_{\RR} y$.
		\item Si $x$ y $y$ son números reales, entonces $x\leq_{\RR} y$ o $y\leq_{\RR} x$.
	\end{enumerate}
\end{theorem}

\begin{proof}
	Consideremos a los números reales $x =\set*{x_{k}}_{k\in\NN}$, $y =\set*{y_{k}}_{k\in\NN}$ y $z =\set*{z_{k}}_{k\in\NN}$.

	\begin{enumerate}
		\item Puesto que para cada $k$ se satisface que $2^{-k + 1} > 0$ se sigue que $x_{k} < x_{k} + 2^{-k + 1}$, es decir, $x\leq_{\RR} x$.
		\item Si $x\leq_{\RR} y$ y $y\leq_{\RR} z$, entonces sabemos que para cada $k$, $x_{k} < y_{k} + 2^{-k + 1}$ y $y_{k} < z_{k} + 2^{-k + 1}$. Consideremos un valor fijo $p > 1$, entonces para cada $k$ se sigue lo siguiente:
		\begin{align*}
			x_{k} &< y_{k} + 2^{-k + 1}\\
			&< y_{k} - y_{k + p} + y_{k + p} + 2^{-k + 1}\\
			&< 2^{-k} + y_{k + p} + 2^{-k + 1}\\
			&< 2^{-k} + z_{k + p} + 2^{-k - p + 1} + 2^{-k + 1}\\
			&< z_{k + p} - z_{k} + z_{k} + 2^{-k} + 2^{-k + 1} + 2^{-k - p + 1}\\
			&< 2^{-k} + z_{k} + 2^{-k} + 2^{-k + 1} + 2^{-k - p + 1}\\
			&= z_{k} + 2^{-k + 1}2^{-p + 1}\\
			&< z_{k} + 2^{-k + 1}
		\end{align*}
		Por lo tanto, $x\leq_{\RR} z$.
		\item Si $x\leq_{\RR} y$ y $y\leq_{\RR} x$, entonces para cada $k$, $x_{k} < y_{k} + 2^{-k + 1}$ y $y_{k} < x_{k} + 2^{-k + 1}$. De la segunda desigualdad obtenemos que $y_{k} - 2^{-k + 1} < x_{k}$, por lo que $y_{k} - 2^{k + 1} < x_{k} < y_{k} + 2^{-k + 1}$, es decir, $\abs*{x_{k} - y_{k}} < 2^{-k + 1}$ para cada $k$. Por lo tanto, $x =_{\RR} y$.
	\end{enumerate}
\end{proof}
